\chapter{Differential Forms and Geometry}

\begin{remarkbox}
Differential forms provide a compact language for integration, Stokes' theorem,
Maxwell theory, Yang--Mills theory, and geometric field theories. This appendix
collects the minimal definitions and identities used in the notes.
\end{remarkbox}

\section{Forms and Exterior Derivative}

A \(p\)-form on a manifold is an antisymmetric covariant tensor field of degree
\(p\). In coordinates,
\[
    \alpha=\frac{1}{p!}\alpha_{\mu_1\cdots\mu_p}
    dx^{\mu_1}\wedge\cdots\wedge dx^{\mu_p}.
\]
A zero-form is an ordinary function. A one-form is
\[
    A=A_\mu dx^\mu.
\]
A two-form is
\[
    F=\frac{1}{2}F_{\mu\nu}dx^\mu\wedge dx^\nu.
\]

The exterior derivative maps \(p\)-forms to \((p+1)\)-forms:
\[
    d:\Omega^p(M)\to\Omega^{p+1}(M).
\]
For a function \(f\),
\[
    df=\partial_\mu f\,dx^\mu.
\]
For a one-form \(A=A_\mu dx^\mu\),
\[
    dA=\frac{1}{2}(\partial_\mu A_\nu-\partial_\nu A_\mu)
    dx^\mu\wedge dx^\nu.
\]
The fundamental identity is
\[
    d^2=0.
\]

\begin{definitionbox}
The exterior derivative is the coordinate-independent derivative on differential
forms. Its defining structural property is \(d^2=0\).
\end{definitionbox}

The identity \(d^2=0\) is the geometric origin of many homogeneous field
equations, including the Bianchi identity \(dF=0\) in electromagnetism when
\(F=dA\).

\section{Wedge Product}

The wedge product combines a \(p\)-form and a \(q\)-form into a \((p+q)\)-form:
\[
    \wedge:\Omega^p(M)\times\Omega^q(M)\to\Omega^{p+q}(M).
\]
It is graded-antisymmetric:
\[
    \alpha\wedge\beta=(-1)^{pq}\beta\wedge\alpha
\]
for \(\alpha\in\Omega^p(M)\) and \(\beta\in\Omega^q(M)\). In particular,
\[
    dx^\mu\wedge dx^\nu=-dx^\nu\wedge dx^\mu.
\]
Therefore
\[
    dx^\mu\wedge dx^\mu=0
\]
with no sum.

The exterior derivative satisfies the graded Leibniz rule:
\[
    d(\alpha\wedge\beta)=d\alpha\wedge\beta+(-1)^p\alpha\wedge d\beta
\]
for \(\alpha\in\Omega^p(M)\).

For Lie-algebra-valued forms, the wedge product also includes multiplication in
the Lie algebra or matrix algebra. Thus for a connection one-form \(A\), the
expression
\[
    A\wedge A
\]
contains both form antisymmetry and noncommuting matrix multiplication.

\section{Hodge Dual}

The Hodge dual maps \(p\)-forms in \(n\) dimensions to \((n-p)\)-forms:
\[
    \star:\Omega^p(M)\to\Omega^{n-p}(M).
\]
It depends on the metric and orientation. In components, for a \(p\)-form
\(\alpha\), the dual has components built using the Levi-Civita tensor and the
metric.

In four-dimensional electromagnetism, the field strength is a two-form
\[
    F=\frac{1}{2}F_{\mu\nu}dx^\mu\wedge dx^\nu,
\]
and its Hodge dual is again a two-form:
\[
    \star F=\frac{1}{2}\widetilde{F}_{\mu\nu}dx^\mu\wedge dx^\nu.
\]
The source-free Maxwell equations can be written compactly as
\[
    dF=0,
    \qquad
    d\star F=0.
\]
With sources, one writes schematically
\[
    d\star F=\star J.
\]

\begin{remarkbox}
Unlike the exterior derivative, the Hodge dual is not purely topological. It
requires a metric. This is why equations involving \(\star\) know about causal
or metric structure.
\end{remarkbox}

\section{Stokes' Theorem}

Stokes' theorem is the central integration theorem for forms. If \(\alpha\) is
an \((n-1)\)-form on an \(n\)-dimensional oriented region \(\Omega\), then
\[
    \int_\Omega d\alpha=\int_{\partial\Omega}\alpha.
\]
This is the common geometric origin of the fundamental theorem of calculus,
Green's theorem, Gauss's theorem, and the ordinary Stokes theorem of vector
calculus.

In variational calculus, integration by parts is a local expression of Stokes'
theorem. For example,
\[
    \int_\Omega d^{d+1}x\,\partial_\mu V^\mu
    =\int_{\partial\Omega}d\Sigma_\mu V^\mu.
\]
Boundary terms in field theory are therefore not artificial. They are required
by Stokes' theorem.

\begin{definitionbox}
Stokes' theorem converts a bulk integral of an exterior derivative into a
boundary integral:
\[
    \int_\Omega d\alpha=\int_{\partial\Omega}\alpha.
\]
\end{definitionbox}

This theorem is also the reason why total derivative terms can affect boundary
charges while leaving local bulk equations unchanged.

\section{Connections and Curvature}

A connection allows one to differentiate fields that carry internal or geometric
indices in a way compatible with local frame changes. In Abelian gauge theory,
the connection is the electromagnetic potential
\[
    A=A_\mu dx^\mu.
\]
The curvature is
\[
    F=dA.
\]

In non-Abelian gauge theory, the connection is Lie-algebra-valued:
\[
    A=A_\mu^aT_a dx^\mu.
\]
The curvature is
\[
    F=dA+gA\wedge A.
\]
In components,
\[
    F_{\mu\nu}=\partial_\mu A_\nu-\partial_\nu A_\mu+g[A_\mu,A_\nu].
\]

The covariant exterior derivative of a Lie-algebra-valued form \(X\) is
schematically
\[
    DX=dX+g[A,X],
\]
with signs depending on the form degree. In components for an adjoint-valued
field,
\[
    D_\mu X=\partial_\mu X+g[A_\mu,X].
\]

\begin{definitionbox}
The curvature of a connection is the field strength. It measures the failure of
covariant derivatives to commute:
\[
    [D_\mu,D_\nu]=gF_{\mu\nu}
\]
when acting on a field in the relevant representation.
\end{definitionbox}

\section{Bianchi Identities}

For electromagnetism,
\[
    F=dA.
\]
Therefore
\[
    dF=d^2A=0.
\]
This is the Abelian Bianchi identity and gives the homogeneous Maxwell equations.

For a non-Abelian connection,
\[
    F=dA+gA\wedge A.
\]
The Bianchi identity is covariant:
\[
    DF=0.
\]
In components,
\[
    D_{[\lambda}F_{\mu\nu]}=0.
\]
This is not an equation of motion. It follows from the definition of curvature.

Bianchi identities also appear in geometry. For the Riemann curvature tensor,
the differential Bianchi identity is
\[
    \nabla_{[\lambda}R_{\mu\nu]\rho}{}^\sigma=0.
\]
After contractions it leads to
\[
    \nabla_\mu G^{\mu\nu}=0,
\]
where \(G^{\mu\nu}\) is the Einstein tensor.

\section{Applications to Maxwell and Yang--Mills Theory}

Maxwell theory in form notation is especially compact. The field strength is
\[
    F=dA.
\]
The homogeneous equations are
\[
    dF=0.
\]
The inhomogeneous equations are
\[
    d\star F=\star J.
\]
The action is
\[
    S[A]=-\frac{1}{2}\int F\wedge\star F
\]
up to conventional normalizations.

Yang--Mills theory replaces \(F=dA\) by
\[
    F=dA+gA\wedge A.
\]
The action is
\[
    S[A]=-\frac{1}{2}\int \langle F\wedge\star F\rangle.
\]
The equation of motion is
\[
    D\star F=\star J,
\]
and the Bianchi identity is
\[
    DF=0.
\]

The form notation makes the analogy between Maxwell and Yang--Mills theory
transparent. The difference is that the non-Abelian theory uses a covariant
exterior derivative and a curvature with a quadratic connection term.

\section{Applications to Gravity}

Gravity can also be written geometrically using connections and curvature. In
ordinary differential geometry, the Levi-Civita connection defines covariant
derivatives compatible with the metric and torsion-free condition. Its curvature
is the Riemann tensor.

In first-order formulations of gravity, one uses a coframe \(e^a\) and a spin
connection \(\omega^{ab}\). The torsion two-form is
\[
    T^a=de^a+\omega^a{}_b\wedge e^b,
\]
and the curvature two-form is
\[
    R^{ab}=d\omega^{ab}+\omega^a{}_c\wedge\omega^{cb}.
\]
The Bianchi identities become
\[
    DT^a=R^a{}_b\wedge e^b,
\]
and
\[
    DR^{ab}=0.
\]

These formulas show that gauge theory and gravity share a geometric language:
connections, curvature, covariant derivatives, Bianchi identities, and boundary
terms. The physical interpretation differs, but the mathematical pattern is the
same.

\begin{summarybox}
Differential forms package many field-theory structures into compact geometric
form. Electromagnetism is \(F=dA\), Yang--Mills theory is
\(F=dA+gA\wedge A\), Bianchi identities are curvature identities, and Stokes'
theorem explains why boundary terms are unavoidable in variational principles
and charge constructions.
\end{summarybox}
